\documentclass[12pt,a4paper,openany,oneside]{book}

% --- Paquetes ---
\usepackage[utf8]{inputenc}
\usepackage[spanish, es-tabla]{babel}
\usepackage{amsmath, amsfonts, amssymb}
\usepackage{graphicx}
\usepackage{geometry}
\usepackage{hyperref}
\usepackage{booktabs}
\usepackage{fancyhdr}
\usepackage{listings}
\usepackage{xcolor}
\usepackage{tikz}
\usepackage{enumitem}
\usepackage{array}
\usetikzlibrary{arrows.meta, positioning, shapes.geometric}

\geometry{margin=2.5cm}

% --- Estilo de Código Python ---
\definecolor{codegreen}{rgb}{0,0.6,0}
\definecolor{codegray}{rgb}{0.5,0.5,0.5}
\definecolor{codepurple}{rgb}{0.58,0,0.82}
\definecolor{backcolour}{rgb}{0.95,0.95,0.92}

\lstset{
    language=Python,
    backgroundcolor=\color{backcolour},   
    commentstyle=\color{codegreen},
    keywordstyle=\color{blue},
    numberstyle=\tiny\color{codegray},
    stringstyle=\color{codepurple},
    basicstyle=\ttfamily\small,
    breaklines=true,
    frame=single,
    showstringspaces=false
}

% --- Cabecera y Pie de Página ---
\pagestyle{fancy}
\fancyhf{}
\renewcommand{\headrulewidth}{0.4pt}
\renewcommand{\footrulewidth}{0.4pt}
\fancyhead[L]{\small Carlos César Sánchez Coronel}
\fancyhead[R]{\small Cálculo y Álgebra Lineal para IA}
\fancyfoot[L]{\small Carlos César Sánchez Coronel}
\fancyfoot[C]{\thepage}

% --- Portada ---
\title{
    \textbf{Sílabo del Curso} \\ 
    \Huge \textbf{Cálculo y Álgebra Lineal} \\
    \Large \textsc{para Inteligencia Artificial} \\[1cm]

}
\author{
    \small Por \\ \vspace{0.2cm}
    \Large \textsc{Carlos César Sánchez Coronel}
}
\date{2026}

\begin{document}

\maketitle
\tableofcontents

% ------------------------------------------------------------------
% PRESENTACIÓN DEL CURSO
% ------------------------------------------------------------------
\chapter*{Presentación del Curso}
\addcontentsline{toc}{chapter}{Presentación del Curso}

Este curso ha sido diseñado para futuros ingenieros en inteligencia artificial. Las matemáticas son el lenguaje de la IA: cada algoritmo de \textit{machine learning}, cada red neuronal, cada modelo generativo se sustenta en conceptos de cálculo diferencial, integral y álgebra lineal. 

A lo largo de 14 semanas, exploraremos estos fundamentos con un enfoque práctico y orientado a aplicaciones reales: desde el gradiente descendente que entrena una red hasta la descomposición en valores singulares que reduce la dimensionalidad de miles de imágenes. Utilizaremos Python como herramienta de cálculo y simulación, con cuadernos de Jupyter diseñados específicamente para este curso.

Cada tema estará conectado explícitamente con cursos futuros: \textbf{Machine Learning, Deep Learning, Visión por Computador, Robótica, Procesamiento de Lenguaje Natural (NLP) y Modelos Generativos (LLMs)}. El objetivo es que al finalizar, el estudiante no solo domine los conceptos, sino que comprenda profundamente \textit{por qué} son la base de la inteligencia artificial moderna.

\vspace{0.5cm}
\textbf{Estructura del curso:}
\begin{itemize}
    \item[$\square$] Bloque I: Fundaciones Matemáticas y Computacionales (Semanas 1-4)
    \item[$\square$] Bloque II: Cálculo Diferencial y el Motor del Aprendizaje (Semanas 5-9)
    \item[$\square$] Bloque III: Álgebra Lineal Profunda: Transformaciones y Descomposiciones (Semanas 10-14)
\end{itemize}

% ------------------------------------------------------------------
% BLOQUE I: FUNDACIONES MATEMÁTICAS Y COMPUTACIONALES (Semanas 1-4)
% ------------------------------------------------------------------
\chapter{Fundaciones Matemáticas y Computacionales (Semanas 1-4)}
\section*{Objetivo del Bloque}
Construir el puente entre el álgebra de secundaria, la programación con Python y los primeros conceptos de IA.

% ---------------- Semana 1 ----------------
\section{Semana 1: El Lenguaje de la IA: Datos, Funciones y Python}

\subsection{Conceptos}
\begin{itemize}
    \item Repaso de conjuntos numéricos: escalares, vectores.
    \item Funciones como ``cajas negras'': dominio, rango, tipos (lineal, polinomial, logarítmica, exponencial).
    \item Lógica proposicional y su relación con el flujo de código.
    \item Notación de sumatorias y productorias.
\end{itemize}

\subsection{Aplicaciones en IA y Conexión con Cursos Futuros}
\begin{itemize}
    \item \textbf{Machine Learning:} Visualización de datos con Matplotlib. ¿Por qué la función de pérdida \textit{Cross-Entropy} usa logaritmos?
    \item \textbf{Deep Learning:} Introducción a funciones de activación (Sigmoide, ReLU) como ejemplos de funciones no lineales. Cada neurona es una función compuesta.
    \item \textbf{NLP:} La base para representar palabras como vectores (embeddings).
\end{itemize}

\subsection{Python}
\begin{itemize}
    \item Primeros pasos con NumPy (creación de \texttt{arrays}) y Matplotlib (gráficas de funciones).
\end{itemize}

% ---------------- Semana 2 ----------------
\section{Semana 2: Tensores: La Estrella del Data Science}

\subsection{Conceptos}
\begin{itemize}
    \item Escalares (0D-tensor), Vectores (1D), Matrices (2D), Tensores (nD).
    \item Transposición de tensores.
    \item Normas de vectores (L1, L2) como medida de magnitud.
\end{itemize}

\subsection{Aplicaciones en IA y Conexión con Cursos Futuros}
\begin{itemize}
    \item \textbf{Visión por Computador (CV):} Representación de datos: un escalar es un precio, un vector es una imagen en escala de grises (aplanada), una matriz es una imagen RGB, un tensor de 4D es un lote de videos (batch, frames, alto, ancho, canales).
    \item \textbf{Machine Learning:} La norma L2 como regularizador (Ridge) y la L1 como regularizador (Lasso) para evitar sobreajuste.
    \item \textbf{NLP:} Los embeddings de palabras son vectores (1D). Los transformers trabajan con tensores de 3D (batch, secuencia, dimensión).
\end{itemize}

\subsection{Python}
\begin{itemize}
    \item Manipulación de tensores con NumPy (\texttt{.shape}, \texttt{.reshape}, \texttt{.T}).
    \item Cálculo de normas con \texttt{np.linalg.norm}.
\end{itemize}

% ---------------- Semana 3 ----------------
\section{Semana 3: Operaciones con Tensores I: Construyendo Algoritmos}

\subsection{Conceptos}
\begin{itemize}
    \item Suma, resta, producto Hadamard (element-wise).
    \item \textit{Broadcasting}.
    \item Reducción (suma, media a lo largo de ejes).
    \item Producto punto (\textit{dot product}).
\end{itemize}

\subsection{Aplicaciones en IA y Conexión con Cursos Futuros}
\begin{itemize}
    \item \textbf{LLMs y Transformers:} El producto punto como medida de similitud (ej. similitud de coseno para \textit{embeddings}). El mecanismo de \textbf{Atención} en modelos como GPT se basa en productos punto entre consultas y claves ($QK^T$).
    \item \textbf{Deep Learning:} \textit{Broadcasting} para normalizar un lote de imágenes (\textit{Batch Normalization}).
    \item \textbf{Machine Learning:} Reducción para calcular la pérdida promedio en un \textit{batch} durante el entrenamiento.
\end{itemize}

\subsection{Python}
\begin{itemize}
    \item Operaciones de arrays con NumPy.
    \item \textit{Broadcasting} en acción.
    \item Implementación de similitud de coseno.
\end{itemize}

% ---------------- Semana 4 ----------------
\section{Semana 4: Sistemas de Ecuaciones y la Base de la Regresión}

\subsection{Conceptos}
\begin{itemize}
    \item Sistemas de ecuaciones lineales. Representación matricial \(Ax = b\).
    \item Resolución por eliminación (sustitución, eliminación gaussiana).
    \item Sistemas sobredeterminados y la solución por mínimos cuadrados.
\end{itemize}

\subsection{Aplicaciones en IA y Conexión con Cursos Futuros}
\begin{itemize}
    \item \textbf{Machine Learning:} La regresión lineal como la reina de los sistemas sobredeterminados. Encontrar la línea que mejor se ajusta a un conjunto de puntos es el primer modelo predictivo.
    \item \textbf{Robótica:} Resolución de sistemas de ecuaciones cinemáticos.
    \item \textbf{Optimización:} Los mínimos cuadrados son la base de muchos problemas de optimización convexa.
\end{itemize}

\subsection{Python}
\begin{itemize}
    \item Resolver sistemas con \texttt{np.linalg.solve}.
    \item Implementar mínimos cuadrados desde cero y comparar con \texttt{sklearn.linear\_model.LinearRegression}.
\end{itemize}

% ------------------------------------------------------------------
% BLOQUE II: CÁLCULO DIFERENCIAL Y EL MOTOR DEL APRENDIZAJE (Semanas 5-9)
% ------------------------------------------------------------------
\chapter{Cálculo Diferencial y el Motor del Aprendizaje (Semanas 5-9)}
\section*{Objetivo del Bloque}
Comprender cómo las máquinas ``aprenden'' ajustando parámetros mediante el cálculo de derivadas, culminando en la implementación manual del descenso de gradiente y su automatización con librerías.

% ---------------- Semana 5 ----------------
\section{Semana 5: Límites y Continuidad}

\subsection{Conceptos}
\begin{itemize}
    \item Concepto intuitivo de límite.
    \item Límites laterales y al infinito.
    \item Continuidad y tipos de discontinuidades.
\end{itemize}

\subsection{Aplicaciones en IA y Conexión con Cursos Futuros}
\begin{itemize}
    \item \textbf{Deep Learning:} La función escalón (discontinua) vs. funciones de activación como Sigmoide o ReLU (continuas a trozos). Preferimos funciones suaves (diferenciables) para el gradiente descendente.
    \item \textbf{Optimización:} El comportamiento asintótico de las funciones de pérdida (ej. cómo la pérdida tiende a cero a medida que el modelo aprende).
\end{itemize}

\subsection{Python}
\begin{itemize}
    \item Cálculo de límites simbólicos con \texttt{SymPy}.
    \item Visualización de continuidad con Matplotlib.
\end{itemize}

% ---------------- Semana 6 ----------------
\section{Semana 6: La Derivada: Midiendo el Cambio}

\subsection{Conceptos}
\begin{itemize}
    \item Definición de derivada como tasa de cambio. Interpretación geométrica (pendiente de la recta tangente).
    \item Reglas de derivación (potencia, constante, suma, producto, cociente).
    \item Derivadas de funciones comunes (exponencial, logaritmo).
\end{itemize}

\subsection{Aplicaciones en IA y Conexión con Cursos Futuros}
\begin{itemize}
    \item \textbf{Machine Learning:} Derivar la función de pérdida de Error Cuadrático Medio (MSE). Entender por qué la derivada nos dice cómo ajustar un peso para minimizar el error.
    \item \textbf{Deep Learning:} La derivada es la base del \textit{backpropagation}.
\end{itemize}

\subsection{Python}
\begin{itemize}
    \item Cálculo simbólico con \texttt{SymPy}.
    \item Derivación numérica y comparación con la analítica.
\end{itemize}

% ---------------- Semana 7 ----------------
\section{Semana 7: Regla de la Cadena y Derivadas de Orden Superior}

\subsection{Conceptos}
\begin{itemize}
    \item Composición de funciones. La regla de la cadena.
    \item Derivadas de orden superior (segunda derivada).
\end{itemize}

\subsection{Aplicaciones en IA y Conexión con Cursos Futuros}
\begin{itemize}
    \item \textbf{Deep Learning:} Una red neuronal es una función gigante compuesta. La regla de la cadena nos permite calcular cómo el error final depende de los pesos de las capas intermedias (Backpropagation).
    \item \textbf{Optimización:} La segunda derivada (Matriz Hessiana) nos informa sobre la curvatura del error y se usa en optimizadores de segundo orden.
\end{itemize}

\subsection{Python}
\begin{itemize}
    \item Ejemplo numérico de backpropagation en una neurona con dos capas.
    \item Cálculo de la segunda derivada con SymPy.
\end{itemize}

% ---------------- Semana 8 ----------------
\section{Semana 8: Derivadas Parciales y el Vector Gradiente}

\subsection{Conceptos}
\begin{itemize}
    \item Derivadas parciales.
    \item El vector gradiente (\(\nabla\)). Interpretación: dirección y magnitud de máximo crecimiento.
    \item Visualización de curvas de nivel y mapas de gradiente.
\end{itemize}

\subsection{Aplicaciones en IA y Conexión con Cursos Futuros}
\begin{itemize}
    \item \textbf{Machine/ Deep Learning:} El gradiente de la función de coste nos indica la dirección de ``cuesta arriba''. Para minimizar, vamos en dirección opuesta (descenso de gradiente). Cada peso tiene su propia derivada parcial.
    \item \textbf{Visión por Computador:} El gradiente se usa en detección de bordes (filtros Sobel, Canny).
\end{itemize}

\subsection{Python}
\begin{itemize}
    \item Cálculo de gradientes de funciones simples con \texttt{SymPy}.
    \item Visualizar el gradiente sobre una superficie de coste 3D (función de pérdida).
\end{itemize}

% ---------------- Semana 9 ----------------
\section{Semana 9: Automatizando el Gradiente: Introducción a PyTorch/TensorFlow y Descenso de Gradiente}

\subsection{Conceptos}
\begin{itemize}
    \item Algoritmo de descenso de gradiente. Tasa de aprendizaje (\textit{learning rate}). Épocas.
    \item Variantes: Batch, Mini-batch, Estocástico (SGD).
    \item Diferenciación automática (\textit{Automatic Differentiation}). Grafos computacionales.
\end{itemize}

\subsection{Aplicaciones en IA y Conexión con Cursos Futuros}
\begin{itemize}
    \item \textbf{Deep Learning:} Cómo las librerías modernas (PyTorch/ TensorFlow) calculan gradientes por nosotros, liberándonos de derivar a mano. Reescribir el entrenamiento de una neurona usando \texttt{torch.autograd} o \texttt{tf.GradientTape}. Esto es exactamente lo que se hace en la práctica para entrenar modelos con millones de parámetros.
    \item \textbf{Todos los cursos futuros:} Todos los optimizadores (Adam, RMSprop) son variantes de esta idea.
\end{itemize}

\subsection{Python}
\begin{itemize}
    \item Implementar descenso de gradiente desde cero con NumPy.
    \item Primeros pasos con PyTorch/TensorFlow: tensores, gradientes automáticos, capas simples.
    \item Comparar la implementación manual con la automática.
\end{itemize}

% ------------------------------------------------------------------
% BLOQUE III: ÁLGEBRA LINEAL PROFUNDA: TRANSFORMACIONES Y DESCOMPOSICIONES (Semanas 10-14)
% ------------------------------------------------------------------
\chapter{Álgebra Lineal Profunda: Transformaciones y Descomposiciones (Semanas 10-14)}
\section*{Objetivo del Bloque}
Dominar las operaciones avanzadas que permiten comprimir, transformar y entender la estructura de los datos, con aplicaciones directas en modelos de ML/DL.

% ---------------- Semana 10 ----------------
\section{Semana 10: Espacios Vectoriales, Cambio de Base y Transformaciones Lineales}

\subsection{Conceptos}
\begin{itemize}
    \item Espacios vectoriales, subespacios.
    \item Combinación lineal, independencia lineal. Bases y dimensión.
    \item Transformaciones lineales (rotación, escalado, cizallamiento) y su representación matricial.
\end{itemize}

\subsection{Aplicaciones en IA y Conexión con Cursos Futuros}
\begin{itemize}
    \item \textbf{Deep Learning:} Las capas de una red neuronal son transformaciones lineales (la parte de pesos \(Wx\)) seguidas de no linealidades. Una red ``deforma'' el espacio de entrada para hacer los datos linealmente separables.
    \item \textbf{NLP:} El cambio de base es análogo a cambiar la representación de un embedding (ej. de one-hot a un espacio semántico denso).
    \item \textbf{Visión por Computador:} Las transformaciones afines (rotación, traslación) se usan en \textit{data augmentation}.
\end{itemize}

\subsection{Python}
\begin{itemize}
    \item Visualizar el efecto de una matriz de transformación sobre un conjunto de puntos en 2D.
    \item Implementar una capa lineal simple con NumPy.
\end{itemize}

% ---------------- Semana 11 ----------------
\section{Semana 11: Determinantes, Inversas y Sistemas Lineales}

\subsection{Conceptos}
\begin{itemize}
    \item Determinante de una matriz (interpretación geométrica como factor de escala de volumen).
    \item Matriz inversa y condición de invertibilidad.
    \item Rango de una matriz, nulidad.
\end{itemize}

\subsection{Aplicaciones en IA y Conexión con Cursos Futuros}
\begin{itemize}
    \item \textbf{Machine Learning:} El determinante de la matriz de covarianza aparece en la función de densidad de la distribución normal multivariante.
    \item \textbf{Robótica:} El determinante del Jacobiano relaciona velocidades en el espacio de articulaciones con velocidades en el espacio de tareas.
    \item \textbf{Redes Neuronales:} El rango de una matriz de pesos indica la cantidad de ``información'' o ``capacidad'' de esa capa (evitar matrices de rango bajo que limiten la expresividad).
\end{itemize}

\subsection{Python}
\begin{itemize}
    \item Cálculo de determinantes, inversas y rango con NumPy (\texttt{np.linalg.det}, \texttt{np.linalg.inv}, \texttt{np.linalg.matrix\_rank}).
    \item Resolver sistemas lineales con inversa y comparar eficiencia.
\end{itemize}

% ---------------- Semana 12 ----------------
\section{Semana 12: Autovalores y Autovectores: Las Direcciones Clave}

\subsection{Conceptos}
\begin{itemize}
    \item Definición de autovalor (\(\lambda\)) y autovector (\(v\)) para una matriz \(A\) (\(Av = \lambda v\)).
    \item Cálculo manual (2x2) y con Python.
    \item Interpretación geométrica: direcciones que no cambian bajo la transformación, solo se escalan.
\end{itemize}

\subsection{Aplicaciones en IA y Conexión con Cursos Futuros}
\begin{itemize}
    \item \textbf{Machine Learning:} \textbf{PCA (Análisis de Componentes Principales):} Los autovectores de la matriz de covarianza son las direcciones de máxima varianza en los datos. Permiten reducir dimensionalidad (de cientos de variables a unas pocas componentes) para visualización o para alimentar modelos más simples.
    \item \textbf{Sistemas de Recomendación:} El algoritmo PageRank de Google es un problema de autovectores (el autovector principal de la matriz de la web).
    \item \textbf{Deep Learning:} Análisis de la matriz de pesos para entender la dinámica del aprendizaje.
\end{itemize}

\subsection{Python}
\begin{itemize}
    \item Calcular autovalores/autovectores con \texttt{np.linalg.eig}.
    \item Aplicarlos para implementar PCA desde cero y visualizar la reducción de dimensionalidad en un dataset como Iris o dígitos.
\end{itemize}

% ---------------- Semana 13 ----------------
\section{Semana 13: La Navaja Suiza: Descomposición en Valores Singulares (SVD)}

\subsection{Conceptos}
\begin{itemize}
    \item Teorema SVD: \(A = U\Sigma V^T\).
    \item Interpretación de las matrices \(U\) (autovectores izquierdos), \(\Sigma\) (valores singulares), \(V^T\) (autovectores derechos) como rotaciones y escalados.
    \item Pseudoinversa de Moore-Penrose.
\end{itemize}

\subsection{Aplicaciones en IA y Conexión con Cursos Futuros}
\begin{itemize}
    \item \textbf{Visión por Computador:} Compresión de imágenes (aproximación de bajo rango). Reducción de ruido.
    \item \textbf{Sistemas de Recomendación:} Filtrado colaborativo (descomponer la matriz de usuarios-productos en factores latentes). La base de Netflix Prize.
    \item \textbf{NLP:} \textit{Latent Semantic Analysis (LSA)} para encontrar temas en documentos.
    \item \textbf{Deep Learning:} La pseudoinversa se usa para resolver mínimos cuadrados de forma analítica y para inicializar ciertos tipos de redes.
\end{itemize}

\subsection{Python}
\begin{itemize}
    \item Aplicar SVD con \texttt{np.linalg.svd} para comprimir una imagen y ver el efecto de diferentes rangos.
    \item Usar la pseudoinversa para resolver un sistema de regresión lineal.
\end{itemize}

% ---------------- Semana 14 ----------------
\section{Semana 14: Integrales y Aplicaciones en Evaluación de Modelos}

\subsection{Conceptos}
\begin{itemize}
    \item Integral indefinida (antiderivada).
    \item Integral definida como área bajo la curva.
    \item Teorema Fundamental del Cálculo.
\end{itemize}

\subsection{Aplicaciones en IA y Conexión con Cursos Futuros}
\begin{itemize}
    \item \textbf{Machine Learning:} \textbf{Curva ROC y AUC (Area Under the Curve):} La integral (área bajo la curva) es una métrica clave para evaluar clasificadores binarios. Cuantifica la capacidad del modelo para distinguir entre clases independientemente del umbral.
    \item \textbf{Probabilidad:} La integral de una función de densidad de probabilidad (PDF) nos da la probabilidad acumulada (CDF). El valor esperado \(E[X] = \int x f(x) dx\) es una integral.
    \item \textbf{Deep Learning (Generativo):} En modelos como \textit{Variational Autoencoders (VAEs)}, se optimiza una cota inferior de la verosimilitud (ELBO) que involucra integrales (expectativas).
\end{itemize}

\subsection{Python}
\begin{itemize}
    \item Calcular integrales definidas con \texttt{scipy.integrate.quad}.
    \item Calcular el área bajo la curva ROC usando integración numérica y comparar con \texttt{sklearn.metrics.roc\_auc\_score} sobre un clasificador entrenado.
\end{itemize}

% ------------------------------------------------------------------
% METODOLOGÍA Y EVALUACIÓN (14 Semanas)
% ------------------------------------------------------------------
\chapter{Metodología y Evaluación}

\section{Estrategias metodológicas}
\begin{itemize}
    \item \textbf{Clases síncronas:} Exposición conceptual con diapositivas, resolución de problemas en pizarra y programación en vivo con Python utilizando cuadernos de Jupyter diseñados para el curso.
    \item \textbf{Trabajo asíncrono:} Lectura de material complementario y resolución de ejercicios prácticos.
    \item \textbf{Aprendizaje activo:} Discusión de aplicaciones reales en IA (como las detalladas en cada semana), resolución de problemas en equipo y pequeños retos de programación.
    \item \textbf{Recursos:} Plataforma virtual, Zoom académico, grabaciones de clases, foros de consulta y cuadernos de Jupyter.
\end{itemize}

\section{Materiales educativos}
\begin{itemize}
    \item Presentaciones interactivas (PPT, videos).
    \item Cuadernos de Jupyter con ejemplos y ejercicios.
    \item Bibliografía: Apuntes del curso, artículos de divulgación, documentación de NumPy/SciPy/PyTorch.
    \item Google Colab para ejecución de código sin instalación.
\end{itemize}

\section{Evaluación}
\begin{table}[h]
\centering
\begin{tabular}{lc}
\toprule
\textbf{Ítem} & \textbf{Ponderación} \\
\midrule
Evaluaciones continuas (EC) & 40\% \\
Examen parcial (EP) - Semana 7 & 30\% \\
Examen final (EF) - Semana 14 & 30\% \\
\midrule
\textbf{Promedio final (PF)} & PF = 0.4 \times EC + 0.3 \times EP + 0.3 \times EF \\
\bottomrule
\end{tabular}
\caption{Sistema de evaluación}
\end{table}

\begin{itemize}
    \item \textbf{Evaluaciones continuas:} Controles de lectura, entrega de ejercicios de programación, participación en foros.
    \item \textbf{Examen parcial:} Cubre los Bloques I y II (Semanas 1-7).
    \item \textbf{Examen final:} Cubre el Bloque III (Semanas 8-14), con énfasis en aplicaciones.
    \item \textbf{Examen sustitutorio:} Semana 15, reemplaza la nota más baja entre parcial y final.
\end{itemize}

% ------------------------------------------------------------------
% BIBLIOGRAFÍA
% ------------------------------------------------------------------
\chapter*{Bibliografía}
\addcontentsline{toc}{chapter}{Bibliografía}

\begin{itemize}
    \item Strang, G. (2019). \textit{Linear Algebra and Learning from Data}. Wellesley-Cambridge Press.
    \item Goodfellow, I., Bengio, Y., Courville, A. (2016). \textit{Deep Learning}. MIT Press.
    \item Deisenroth, M. P., Faisal, A. A., Ong, C. S. (2020). \textit{Mathematics for Machine Learning}. Cambridge University Press.
    \item Documentación oficial de NumPy, SciPy, Matplotlib, PyTorch, TensorFlow, Scikit-learn.
\end{itemize}

\vspace{1cm}
\begin{flushright}
\textit{Elaborado por:} Mg. Carlos César Sánchez Coronel \\
\textit{Aprobado por:} Dirección de Ingeniería de Inteligencia Artificial \\
\textit{Fecha:} 29/03/2024
\end{flushright}

\end{document}